% =============================================================================
%  slushpile resume template
%
%  Build:  latexmk -xelatex resume.tex && latexmk -c
%  Needs:  XeLaTeX (for fontspec). Fonts degrade gracefully -- see FONTS below.
%
%    Fedora  sudo dnf install -y texlive-xetex texlive-fontspec texlive-microtype \
%                                latexmk dejavu-fonts-all
%    Debian  sudo apt install texlive-xetex texlive-fonts-extra fonts-dejavu latexmk
%    macOS   brew install --cask mactex-no-gui
%
%  Why LaTeX at all: an ATS reads the extracted text layer, and a text-based PDF
%  from XeLaTeX extracts cleanly and deterministically. Word and Google Docs
%  exports vary by exporter version. If you would rather not use LaTeX, write the
%  resume in Markdown -- every skill in this pipeline works on extracted text and
%  none of them require .tex. Only this template does.
%
%  Check the extraction before you send anything:
%
%      pdftotext resume.pdf - | less
%
%  That text is what an ATS parses and what every review agent in this pipeline
%  is handed. If it reads badly, fix the document rather than the review.
%
%  ---------------------------------------------------------------------------
%  ATS RULES BAKED INTO THIS LAYOUT. Do not undo them.
%  ---------------------------------------------------------------------------
%   - Single column. Multi-column layouts get read in the wrong order.
%   - No tables for content. No text boxes. No icons carrying meaning.
%   - Contact details in the body, never in \header or \footer. Roughly a
%     quarter of ATS drop header and footer content entirely.
%   - Standard section names. "Work Experience", not "Where I've Been".
%   - Month and year on every role. Year-only dates make tenure ambiguous, and
%     an ambiguous tenure gets resolved against you.
%   - Every link is also readable as text where it matters. A recruiter reading
%     a printout cannot click \href{}{}, so the visible text is the URL itself.
%
%  ---------------------------------------------------------------------------
%  FONTS
%  ---------------------------------------------------------------------------
%  The house style is Public Sans (body) with IBM Plex Mono (eyebrows, dates).
%  Both are OFL-licensed and both are vendored with this plugin, so installing
%  them is one command and needs no network:
%
%      python3 scripts/install_fonts.py
%
%  Confirm with:
%
%      fc-list | grep -iE 'public sans|plex mono'
%
%  Neither ships with TeX Live, and a missing font is a hard XeLaTeX error
%  rather than a warning: the document does not look wrong, it does not build.
%  So the preamble names each family through \IfFontExistsTF and falls back to
%  DejaVu. Skipping the install is a supported outcome. The documents build and
%  they are set in something else.
%
%  The families are named here rather than pointed at by path on purpose. This
%  file gets copied into a per-role folder in your workspace, at a depth the
%  plugin does not control, so a relative Path= would resolve to nothing and an
%  absolute one would name a checkout that the next plugin update replaces.
%
%  ---------------------------------------------------------------------------
%  HOUSE STYLE
%  ---------------------------------------------------------------------------
%  No em dashes anywhere in the body: comma, colon, semicolon, or parentheses
%  instead. En dashes for ranges only. The reason is not typographic taste. A
%  document dense with em dashes reads as model output to a fatigued screener,
%  and they bin it before they could tell you why.
%
%  Placeholders are UPPERCASE so the onboard skill's placeholder check catches
%  anything you forgot. Leave them uppercase until you replace them.
% =============================================================================

\documentclass[10pt,letterpaper]{article}

\usepackage[top=0.4in,bottom=0.35in,left=0.55in,right=0.55in]{geometry}
\usepackage{titlesec}
\usepackage{enumitem}
\usepackage{fontspec}
\usepackage{xcolor}
\usepackage{textcomp}
\usepackage{microtype}
\usepackage{hyperref}

% ---- palette ----------------------------------------------------------------
% Graphite for structure, one copper accent for emphasis. One accent, used
% sparingly, reads as designed. Three accents read as a template someone found.
\definecolor{base900}{HTML}{1A1B1E}   % section titles, role titles, the name
\definecolor{ink}{HTML}{181B20}       % body
\definecolor{muted}{HTML}{43505D}     % employer line, contact line
\definecolor{faint}{HTML}{6B7480}     % eyebrows, dates
\definecolor{accent}{HTML}{B05B33}    % rules and bullets
\definecolor{accentink}{HTML}{7C3E20} % links, darkened so they stay legible

\hypersetup{
    colorlinks=true, urlcolor=accentink, linkcolor=accentink,
    pdftitle={NAME -- ROLE}, pdfauthor={NAME}
}

% ---- fonts, with a fallback so this builds on a bare TeX Live ---------------
\IfFontExistsTF{Public Sans}{%
    \setmainfont{Public Sans}[
      UprightFont={* Regular}, BoldFont={* Bold},
      ItalicFont={* Italic}, BoldItalicFont={* Bold Italic}]
    \newfontfamily\heavy{Public Sans}[UprightFont={* ExtraBold}]
}{%
    \setmainfont{DejaVu Sans}
    \newfontfamily\heavy{DejaVu Sans}[UprightFont={* Bold}]
}

\IfFontExistsTF{IBM Plex Mono}{%
    \setmonofont{IBM Plex Mono}[Scale=0.84, UprightFont={* Regular}, BoldFont={* SemiBold}]
}{%
    \setmonofont{DejaVu Sans Mono}[Scale=0.78]
}

\color{ink}
\linespread{0.98}
\setlength{\parindent}{0pt}
\setlength{\parskip}{0pt}
\emergencystretch=3em
\hyphenpenalty=400

% ---- eyebrow: letterspaced uppercase mono, for tags and credential lines ----
\newcommand{\eyebrow}[1]{{\ttfamily\scriptsize\color{faint}\addfontfeatures{LetterSpace=4.0}\MakeUppercase{#1}}}
\newcommand{\dateline}[1]{{\ttfamily\footnotesize\color{faint}#1}}

% ---- section: heavy graphite title over a copper rule ----------------------
\titleformat{\section}
    {\heavy\large\color{base900}}{}{0em}{}[{\color{accent}\titlerule[1.4pt]}]
\titlespacing*{\section}{0pt}{7pt}{2pt}

% ---- \resumesubheading{Employer}{Location}{Title}{Dates} -------------------
\newcommand{\resumesubheading}[4]{%
    \noindent{\bfseries\color{base900}#1} \hfill \dateline{#2} \\
    {\itshape\color{muted}#3} \hfill \dateline{#4}}

% ---- \work{Headline}{EYEBROW TAG}{Body} ------------------------------------
% For public or portfolio work. The headline is a claim, not a project name:
% "A pipeline that refuses to publish a number it cannot defend" gets read.
% "kernel-cve-watch" gets skipped by anyone who does not already know it.
%
% NOTE: \sep is used inside \eyebrow, which applies \MakeUppercase. A \color{}
% argument would be uppercased into an undefined color name, so this stays plain.
\newcommand{\sep}{\,\textbullet\,}
\newcommand{\work}[3]{%
    \noindent{\bfseries\color{base900}#1} \hfill \eyebrow{#2}\\[1pt]
    \noindent #3\par\vspace{3pt}}

\setlist[itemize]{leftmargin=1.2em, label={\small\color{accent}\textbullet}}

\pagestyle{empty}

\begin{document}

% ============================== HEADER ==============================
% Contact details live here, in the body. Never in a \header. See the top note.
%
% Line 2 is a positioning line, not a job title. It states what you do and the
% condition you do it under, and it is rewritten per application. "Senior
% Engineer" tells a screener nothing they cannot read from your last role.
%
% The eyebrow line carries the hard credentials a screener would otherwise have
% to hunt for: work authorization, active certifications, years in discipline,
% licences. Delete what does not apply rather than softening it. Every one of
% these comes out of preferences.yaml and profile.md, never invented here.

{\heavy\fontsize{21}{23}\selectfont\color{base900}NAME}\\[2pt]
{\bfseries\color{accentink}POSITIONING LINE: What You Do, And The Condition You Do It Under}\\[3pt]
{\ttfamily\footnotesize\color{muted}CITY, STATE \ \textbullet\ PHONE \ \textbullet\ EMAIL \ \textbullet\ \href{https://github.com/USER}{github.com/USER} \ \textbullet\ \href{https://linkedin.com/in/USER}{linkedin.com/in/USER}}\\[4pt]
\eyebrow{WORK AUTHORIZATION \textbullet\ CERTIFICATIONS \textbullet\ YEARS IN DISCIPLINE}

\vspace{4pt}
{\color{accent}\hrule height 1.4pt}
\vspace{5pt}

% ============================== SUMMARY ==============================
% Two short paragraphs, both rewritten per application.
%
% Do not paste phrases from the job description here. A screener who reads their
% own posting back at them stops reading, and the ATS gains nothing the Skills
% section is not already giving it.
%
% Paragraph one: lead with the strongest match to the core function. Name the
% years in discipline directly rather than defending their absence. End on what
% makes you unusual rather than what makes you adequate.

\noindent SUMMARY. Open on the thing you own rather than the thing you are
familiar with. Put one checkable number in the first three lines: a scale, a
volume, a population that acts on your output. Close on the combination the rest
of the queue does not have.

\vspace{4pt}

% Paragraph two is optional and it is the strongest move in this template when
% you are making a stretch application. State the boundary of what you have done,
% in your own words, before a screener infers a worse version of it.
%
% Naming a gap yourself converts it from a discovered misrepresentation into a
% disclosed limit, and it buys the rest of the document credibility it cannot
% otherwise earn from a stranger. Do this only where the gap is load-bearing for
% the role. Volunteering an irrelevant weakness is not honesty, it is noise.
%
% Delete this paragraph outright when the role has no such gap.

\noindent Underneath that is YEARS doing ADJACENT WORK, including THE ONE
ACHIEVEMENT THAT TRANSFERS. What I have not done: THE NAMED BOUNDARY, stated
plainly and without hedging. THE HONEST VERSION OF WHAT I DO HAVE, and I would
be learning YOUR SPECIFIC TOOL.

% ============================== PUBLIC WORK ==============================
% Move this section below Work Experience when the role weights tenure over
% public technical credibility. Section order is a per-application decision, and
% whichever section proves your thesis goes first.
%
% Two rules for this section specifically:
%
%   1. Everything here must be checkable in thirty seconds by a stranger who
%      will not call a reference. That is the entire reason it outranks a
%      Projects section: verifiable evidence is disproportionately valuable in a
%      stretch application.
%   2. Deep-link where the on-axis artifact is not the repository root. Landing
%      a reader on a README that buries the relevant part costs you the click.

\section{Public Work}

\work{A headline that states what the work proves, not what the project is called}{DOMAIN \sep LICENSE}{%
\href{https://github.com/USER/PROJECT}{github.com/USER/PROJECT}. What it does, who
uses it, and the number that makes it real. Then the mechanism: the specific
decision you made and the failure it prevents. \textbf{Bold the load-bearing
clause} so a skimming reader gets the argument without reading the block, and
use it at most twice per entry. Bold applied to five clauses is the same as bold
applied to none.}

\work{Second entry}{DOMAIN \sep TAG}{%
\href{https://github.com/USER/PROJECT}{github.com/USER/PROJECT}. Ordering in this
section is an argument. If one project caused another, put them adjacent and in
that order, because the causality is the point a list of projects cannot make.}

% ============================== EXPERIENCE ==============================
% Standard heading. Reverse chronological. Month and year, always.

\section{Work Experience}

\resumesubheading{COMPANY}{City, State}{JOB TITLE}{Month YYYY -- Month YYYY}

\vspace{2pt}

% An orienting line under the employer, in italic. A reader who has never heard
% of this company cannot size anything you did there without it: what the company
% does, how big, under what constraint. One sentence. Delete for employers the
% reader certainly knows.
\noindent\textit{\small What this company does, at what scale, under what
constraint that shaped your work.}

\vspace{2pt}

\begin{itemize}[nosep, topsep=2pt, itemsep=3pt]
    \item \textbf{Every bullet passes the ``So What?'' test: baseline, then
          action, then significance.} ``Cut deployment time 60\%'' is a claim.
          ``Cut deployment time from 45 minutes to 18, unblocking daily releases
          for six teams'' is evidence. The baseline is the part people drop and
          it is the part that makes the number mean anything.
    \item \textbf{Lead with what you decided}, not what you participated in.
          The gap between a senior resume and a mid-level one is almost entirely
          here, and most resumes never capture it.
    \item Put scope in the sentence. ``Ran a program'' and ``Ran a program across
          six teams in three orgs against a \$4M budget'' are read as different
          levels by the same reader.
    \item Vary the length deliberately. Bullets all within five words of each
          other is one of the loudest AI tells a fatigued reader picks up, and
          they pick it up before consciously noticing why. This bullet is short.
\end{itemize}

\vspace{2pt}

\resumesubheading{COMPANY}{City, State}{JOB TITLE}{Month YYYY -- Month YYYY}

\begin{itemize}[nosep, topsep=2pt, itemsep=3pt]
    \item
\end{itemize}

% An early role can collapse to a single line once it stops carrying weight.
% Deleting it entirely creates a gap in the dates, and an unexplained gap gets
% resolved against you by a reader who has other resumes to get through.
\vspace{2pt}
\resumesubheading{COMPANY}{City, State}{EARLY JOB TITLE}{Month YYYY -- Month YYYY}\\[3pt]
\noindent{\small One line naming the ground-level skills, so the date range is
accounted for without spending space defending it.}

% If you have independent, contract, or side work that is real, it belongs here
% with dates, not in a Projects section. An undated Projects entry does not count
% toward years-of-experience in an ATS, and for anyone making a career transition
% that is usually the number that decides the screen.

% ============================== SKILLS ==============================
% Keywords carry more weight here than anywhere else in the document. Three
% rules:
%
%   1. Everything here must be defensible in an interview. The requirements
%      analyst in this pipeline checks each entry against your Work Experience
%      and flags anything unsupported, so an inflated claim gets caught in the
%      review rather than in the interview. Cheaper to be honest now.
%   2. Group them. A flat 40-item comma list reads as keyword stuffing to a
%      human and gains nothing with the machine.
%   3. Put the group the role is actually hiring for first, and expand it. The
%      other groups exist to establish range, and range is worth less than depth
%      on the one axis being screened.

\section{Skills}

{\footnotesize
\noindent\textbf{Category the role is hiring for:} the specific things, named
the way the posting names them, each one traceable to a bullet above\par\vspace{1pt}
\noindent\textbf{Category:} item, item, item\par\vspace{1pt}
\noindent\textbf{Category:} item, item, item\par}

% ============================== EDUCATION ==============================
% A combined "Education \& Certifications" heading parses unreliably in some ATS.
% Keep them separate when both are substantial. Combine them, as below, only when
% neither carries much weight for this role and the space is better spent above.

\section{Certifications \& Education}

{\footnotesize\noindent \textbf{CERTIFICATION} (status, number if it is
verifiable), ISSUING BODY, YYYY. DEGREE, INSTITUTION, YYYY.\par}

\end{document}
